\section{Introduction}
\label{sec:intro}

Between April 2009 and January 2011 the author used \emph{Derive for
Windows}, the computer-algebra system by Texas Instruments, as a notebook: one
writes an expression, asks the program to simplify it, solve it,
differentiate it or take its limit, changes a parameter and looks at what
happens. Twenty-three \texttt{.dfw} files with $792$ expressions survive. It
is experimental mathematics in the sense of~\cite{borwein2008}: rarely a
proof, often a glimpsed regularity, sometimes a mistake.

Derive was withdrawn from the market in 2007 and is no longer maintained;
its files cannot be opened by current mathematical software, but they store
the expressions as text. The recovery, which
consists of decoding the format, translating the expressions into
mathematical notation and checking them numerically, is described in
Appendix~\ref{app:recupero}. The complete result is an integral edition of
more than two hundred pages, in Italian, in which every expression is
reproduced and commented on.

This note is something else: a selection by theme, written for a
mathematical reader. Each section says what the notebook does, states cleanly
what is true in it, proves it or sketches the proof, places it in the
literature and says what remains open. No result is presented as new. The
only proof we have not found in the literature, the one in
Section~\ref{sec:fattoriali}, is elementary.

In summary:
\begin{itemize}[leftmargin=1.4em,itemsep=1pt]
  \item Section~\ref{sec:semiangolo} reconstructs the idea that runs through
        five notebooks, writing the exponential with the half-angle tangent,
        and follows it to Euler, Lambert, Pad\'e and the bilinear
        transformation of digital filters;
  \item Section~\ref{sec:fattoriali} explains a regularity between
        factorials that lasts four steps and proves that it does not recur;
  \item Sections~\ref{sec:media}, \ref{sec:cosinc}, \ref{sec:gradini}
        and~\ref{sec:e} trace four more notebooks back to known facts: the
        mean value theorem for the Helmholtz equation, the Hilbert transform,
        Maslov dequantization, the removal of the odd terms of an error;
  \item Section~\ref{sec:delta} discusses a correct computation with an
        interpretation that does not hold up, and leaves the question open;
  \item Section~\ref{sec:breve} summarizes the other notebooks,
        Section~\ref{sec:refusi} the mistakes.
\end{itemize}

\paragraph{How to read the references.}
The notation \riga{factor.dfw}{4} denotes expression number~4 of the file
\texttt{factor.dfw}, in Derive's numbering; \righe{gauss.dfw}{1}{5} denotes a
range of expressions. The integral edition reproduces all of them with the
same numbering.

\paragraph{Verification.}
Every numerical value quoted in the note is recomputed by a script included in
the repository, \texttt{verifica\_nota.py}, with \texttt{mpmath} at sixty
digits; a second script checks that every decimal in the text is among those
verified, and that the Italian and English versions report the same numbers,
the same references and the same citations.
